% Pursuing Stacks (cotes 134-2 to 134-8): archive cote 134-3, site page 41 (PDF page 42), typescript page 41.
% typescript. Draft compared: ps2.tex:2009-2075.
% AI-transcribed by Claude Opus 5.5 (claude-opus-5-5, adaptive thinking, effort high), Message Batches API, batch ps1 (msgbatch_01Rv5hCaNTtUt4JHx1yqoHPL), 2026-11-12; image review recorded in records/CHECK_DECISIONS_134_3_4.tsv.
\scanpage{134-3/41}
\mspage{\hand{41}}

a "crible"\inserted{[or ``sieve'']} in $A_{/a}$ , namely for an arrow $y\hand{\to} x$ in $A_{/a}$ , if the target x is
in the subcategory, so is the source y. This being so, we define \over{a}{t}he
functor
\[ f' = h(f,u) \ : \ A_{/a}=C \hand{\to} B \]
by the conditions that
\[ f'\hand{\big|}(A_{/a_u} = C_u) \ = f\hand{\big|}C_u \]
\[ f'\hand{\big|}(C\hand{\setminus} C_u) \ = \text{constant functor with value } e_B\hand{\unsure{;}} \]
(where $C\hand{\setminus} C_u$ denotes the obvious \uline{full} subcategory of C , complementary to
$C_u$ )\hand{.} \over{t}{T}his defines f' uniquely, \xout{the} on the objects first, and on the
arrows too because the only arrows left \xout{where} in C where we got still to
define f' are arrows $x \hand{\to} y$ with x \over{\unsure{o}}{i}n $C_u$ and y in $C\hand{\setminus} C_u$ (because $C_u$ is a
crible), but then $f'(y) = e_B$ and we have no choice for $f'(x)\hand{\to} f'(y)$\hand{!}
It\hand{'}s trivial checking that this way we get indeed a \uline{functor}\inserted{$f'$:} $C\hand{\to} B$, thus
the map h(a) is defined - and that this map is functorial with respect to
a, i.e. comes from a map h: $F\mathrm{x} I \hand{\to} F$ as we wanted. T\over{\unsure{n}}{h}e commutativity
of (D) is easily checked\over{\unsure{,}}{:} for the left triangle i.e. that the compositum
$F \hand{\simeq} F\mathrm{x} e_o \hand{\to} F\mathrm{x} I \overset{h}{\hand{\longrightarrow}} F$ is the identity, it comes from the fact that if u
factors through $e_o$ then $C_u = C$ hence $f' = f$ ; for the right triangle, it
comes from the fact that if u factors through $e_l$ th\over{a}{e}n $C_u = \emptyset$ (here we use
the assumption $e_o\hand{\cap}e_l = \emptyset_{A\hand{^\wedge}}$), and hence f' is the constant functor $C\hand{\to} B$
with value $e_B$ . This settles b).

We have still to check a), namely that for any $F\hand{\in} \mathrm{Ob}\, A^{\hand{\wedge}}$, the inclusi-%
on\deleted{s} of \inserted{the objects} $F\mathrm{x} e_i$\hand{\unsure{,}} into $F\mathrm{x} I$ are weak equivalences, or what amounts to the same,
that the projection
\[ F\mathrm{x} I \hand{\longrightarrow} F \]
is a weak equivalence - this will be true in fact for any object I of $A^{\hand{\wedge}}$
which is aspheric. I\over{n}{n}deed, we will prove the stronger result that
$F\mathrm{x} I \hand{\to} F$ is \uline{aspheric}, i.e. that the functor
\[ A_{/F\mathrm{x} I} \hand{\to} A_{/F} \]
is aspheric\hand{\unsure{.}} \over{w}{W}e use for this the criterion of p. 38 , which here transla-%
tes into the condition that for any a \inserted[typed]{in A} (such that we got an $a\hand{\to} F$ i.e. \over{\unsure{f}}{s}uch
that $F(a) \hand{\neq} \emptyset$ , but never mind), the category $A_{/a\mathrm{x} I}$ is aspheric, i.e.
the element $a\mathrm{x} I$ of A\hand{$^\wedge$}is asperic. Again, \uline{as} \over{a}{I} \uline{is aspheric}, we are redu-%
ce\over{a}{d} to ch\over{\unsure{a}}{e}cking that $a\mathrm{x} I \hand{\to} I$ is aspheric , which by the same argument
(with F,I replaced by I , a) boils down to the condition that $\over{\unsure{a}}{A}_{/a\mathrm{x} b}$ is
aspheric for any b in \over{a}{A}. Now this is just condition (T \over{\unsure{1}}{2}), we are through\hand{!}
